% Part: lambda-calculus
% Chapter: lambda-definability
% Section: fixpoints

\documentclass[../../../include/open-logic-section]{subfiles}

\begin{document}

\olfileid{lam}{ldf}{fp}

\olsection{النقاط الثابتة}

نريد، مثلًا، أن نعرّف دالة المضروب تعريفًا عوديًا بحد~$\fn{Fac}$
يحقق
\[
\fn{Fac} \ident \lambd[n][\fn{IsZero}\, n\, \num 1 (\fn{Mult}\, n (\fn{Fac}(\fn{Pred} \, n)))]
\]
فالمراد أن يكون مضروب~$n$ مساويًا لـ~$1$ عند~$n = 0$، ومساويًا
لحاصل ضرب~$n$ في مضروب~$n-1$ فيما عدا ذلك. غير أن هذه العبارة لا
تصلح بذاتها تعريفًا لحد~$\fn{Fac}$؛ لأن~$\fn{Fac}$ نفسه قد عاد
في الطرف الأيمن، والتعريف بالإحالة إلى الذات ليس من قواعد حساب لامبدا.
أما إذا كتبنا، مثلًا،
\[
\fn{Mult} \ident \lambd[ab][a (\fn{Add}\, b) \num 0]
\]
فليس في الطرف الأيمن إلا حدود سبقت معرفتها، ومنها~$\fn{Add}$؛ فيمكن
أن نحل تعريف~$\fn{Add}$ محله، فنحصل على حد لامبدا خالص لـ~$\fn{Mult}$.
وإن اشتمل تعريف~$\fn{Add}$ نفسه على اسم معرف، كما يقع في
$\fn{Add}'$، كررنا الإحلال حتى لا يبقى إلا تركيب لامبدا الخالص.

ولا ينتهي هذا المسلك إذا كان التعريف عوديًا كتعريف~$\fn{Fac}$.
فإحلال العبارة المعرفة لـ~$\fn{Fac}$ محل ورود~$\fn{Fac}$ في الطرف الأيمن
يعطي:
\begin{align*}
  \fn{Fac} & \ident
  \lambd[n][\fn{IsZero}\, n\, \num{1}\,
    (\fn{Mult}\, n
      ((\lambd[n][\fn{IsZero}\, n\, \num{1}\,
        (\fn{Mult}\, n\, (\fn{Fac}(\fn{Pred}\, n)))])
      (\fn{Pred}\, n)))]
\end{align*}
ولا يزال~$\fn{Fac}$ حاضرًا في الطرف الأيمن. وكلما أعدنا الإحلال طال
الحد وبقي الاسم، فلا تبلغ العملية حد لامبدا خالصًا. فهذه الحيلة لا
تعرّف المضروب، ولا سائر الدوال العودية.

لكن في المعادلة العودية دلالة نافعة. فلنفرض أن~$f$ حد يمثل دالة
المضروب، ولنضع
\[
\fn{Fac}' \ident \lambd[g][\lambd[n][\fn{IsZero} \, n \, \num 1\, (\fn{Mult} \, n\, (g (\fn{Pred} n)))]]
\]
وهذا الحد، عند تطبيقه في~$f$، أي~$\fn{Fac}'\,f$، يمثل دالة المضروب
أيضًا. فإذا نظرنا إلى~$\fn{Fac}'$ على أنها دالة تأخذ دالة وتردّ
دالة، كانت قيمة~$\fn{Fac'}\, f$ هي~$f$ نفسها، متى كانت~$f$ دالة
المضروب. وتسمى الدالة~$f$ التي تحقق~$\fn{Fac}' \, f \equal[\beta] f$
\emph{نقطة ثابتة} لـ~$\fn{Fac}'$. ومن ثم فدالة المضروب نقطة ثابتة
لـ~$\fn{Fac}'$.

وفي حساب لامبدا حدود تستخرج لحد معطى نقطة ثابتة؛ وبها نحول حدًا
كـ~$\fn{Fac}'$ إلى تعريف فعلي للمضروب.

\begin{defn}
  \ollabel{defn:Turing-Y}
  \emph{مركّب~Y} هو الحد:
  \[
  Y \ident (\lambd[ux][x(uux)])(\lambd[ux][x(uux)]).
  \]
\end{defn}

\begin{thm}
  للحد~$Y$ الخاصية~$Yg \red g(Yg)$ لأي حد~$g$؛ ولذلك يكون~$Yg$
  نقطة ثابتة لـ~$g$.
\end{thm}

\begin{proof}
  رمز إلى~$(\lambd[ux][x(uux)])$ بـ~$U$ اختصارًا، فيكون~$Y \ident UU$.
  وحينئذ
  \begin{align*}
    Y g &\ident (\lambd[ux][x(uux)])U\,g \\
    &\red (\lambd[x][x(UUx)])g\\
    & \red g(UUg) \ident g(Yg).
  \end{align*}
  فالحدان~$g(Yg)$ و$Yg$ يختزلان كلاهما إلى~$g(Yg)$؛ ومن ثم
  $g(Yg) \equal[\beta] Yg$، وهو عين كون~$Yg$ نقطة ثابتة لـ~$g$.
\end{proof}

ولما كان~$Yg$ موضع اختزال، جاز أن نسلك به اختزالًا لا ينتهي:
\begin{align*}
  Y g &\red g (Y g) \\
    &\red g (g (Y g)) \\
    &\red g(g (g (Y g)))\\
    &\ldots
\end{align*}
فيجوز أن نتصور~$Yg$ كأنه~$g$ مطبقة في ذاتها عددًا لا نهاية له من
المرات. وإذا طبقنا~$g$ فيه مرة أخرى، لم نزد---بهذا التصور---شيئًا؛
فتطبيق~$g$ في~$g$ المطبقة عددًا لا نهاية له من المرات في~$Yg$
يظل تطبيقًا لـ~$g$ في~$Yg$ عددًا لا نهاية له من المرات.

وينبغي التمييز هنا بين مسالك الاختزال. فمتتالية اختزال~$\beta$ السابقة، المبدوءة
بـ~$Yg$، لا تنتهي. وإذا طبقنا~$Yg$ في حد ما، كان من مسالك اختزال
$(Yg)N$ المرور بالحدود~$(g(Yg))N$ ثم~$(g(g(Yg)))N$، وهكذا بلا نهاية.
ومع ذلك قد يوجد مسلك \emph{آخر} ينتهي إلى صورة سوية.

ولنر ذلك في المضروب. عرّف~$\fn{Fac}$ بأنه~$Y\, \fn{Fac}'$، أي نقطة
ثابتة لـ~$\fn{Fac'}$. عندئذ:
\begin{align*}
  \fn{Fac}\,\num 3
  &\red Y\, \fn{Fac}' \, \num 3 \\
  &\red \fn{Fac}' (Y \,\fn{Fac}') \, \num 3 \\
  &\ident (\lambd[x][\lambd[n][\fn{IsZero} \, n\, \num 1\,
      (\fn{Mult}\, n\, (x (\fn{Pred}\, n)))]]) \, \fn{Fac} \, \num 3\\
  & \red \fn{IsZero} \, \num 3\, \num 1\,
      (\fn{Mult}\, \num 3\, (\fn{Fac} (\fn{Pred}\, \num 3))) \\
  &\red  \fn{Mult}\, \num 3 \, (\fn{Fac} \, \num 2).
  \intertext{وبالمثل،}
  \fn{Fac}\,\num 2
  &\red  \fn{Mult}\, \num 2 \, (\fn{Fac} \, \num 1) \\
  \fn{Fac}\,\num 1
  &\red  \fn{Mult}\, \num 1 \, (\fn{Fac} \, \num 0)
  \intertext{لكن}
  \fn{Fac}\,\num 0
  &\red \fn{Fac}' (Y \,\fn{Fac}') \, \num 0 \\
  &\ident (\lambd[x][\lambd[n][\fn{IsZero} \, n\, \num 1\,
      (\fn{Mult}\, n\, (x (\fn{Pred}\, n)))]]) \, \fn{Fac} \, \num 0\\
  & \red \fn{IsZero} \, \num 0\, \num 1\,
      (\fn{Mult}\, \num 0\, (\fn{Fac} (\fn{Pred}\, \num 0))).\\
  &\red \num 1.
  \intertext{ومن ثم، بجمع ذلك،}
  \fn{Fac}\,\num 3
  &\red  \fn{Mult}\, \num 3 \,
  (\fn{Mult} \, \num 2\, (\fn{Mult} \, \num 1\, \num 1)).
\end{align*}

وهذا البناء لا يختص بـ~$\fn{Fac'}$، بل يجري في كل تعريف عودي.
فإذا كانت لنا المعادلة
\begin{align*}
g\,x_1\dots x_n & \equal[\beta] N
\intertext{حيث يجوز أن يحتوي~$N$ على~$g$ و$x_1$، \dots،~$x_n$.
وعندئذ يوجد دائمًا حد~$G \ident (Y \lambd[g][\lambd[x_1\dots x_n][N]])$
بحيث}
G\,x_1 \dots x_n & \equal[\beta] \Subst{N}{G}{g}.
\intertext{إذ بحسب مبرهنة النقطة الثابتة،}
G \ident (Y \lambd[g][\lambd[x_1\dots x_n][N]]) & \red \lambd[g][\lambd[x_1\dots x_n][N]](Y \lambd[g][\lambd[x_1\dots x_n][N]])\\
& \ident (\lambd[g][\lambd[x_1\dots x_n][N]])\,G
\intertext{وبالتالي}
G\,x_1 \dots x_n
& \red (\lambd[g][\lambd[x_1\dots x_n][N]])\,G\,x_1\dots x_n\\
& \red (\lambd[x_1\dots x_n][\Subst{N}{G}{g}])\,x_1\dots x_n\\
  & \red \Subst{N}{G}{g}.
\end{align*}
وهذا عين السلوك الذي نصت عليه المعادلة العودية.

ومركّب~$Y$ في \olref{defn:Turing-Y} من وضع آلان تورنغ. ولألونزو
تشيرش صيغة أخرى نرمز إليها بـ~$Y_C$:
\[
Y_C \ident \lambd[g][(\lambd[x][g(xx)])(\lambd[x][g(xx)])].
\]
ومركّب تشيرش أضعف قليلًا من مركّب تورنغ: فهو يحقق
$Y_Cg \equal[\beta] g(Y_Cg)$، ولا يحقق~$Y_Cg \bred g(Y_Cg)$.
وليكن~$V$ الحد~$\lambd[x][g(xx)]$، حتى يكون~$Y_C \ident \lambd[g][VV]$.
فنجد
\begin{align*}
  VV & \ident (\lambd[x][g(xx)])V \red g(VV)
  \text{ ومن ثم}\\
  Y_Cg & \ident (\lambd[g][VV])g \red VV \red g(VV) \text{، وكذلك}\\
  g(Y_Cg) & \ident g((\lambd[g][VV])g) \red g(VV).
\end{align*}
فالحدان~$Y_Cg$ و$g(Y_Cg)$ يردان إلى حد مشترك هو~$g(VV)$؛ ولذلك
$Y_Cg \equal[\beta] g(Y_Cg)$. وهذه المساواة تكفي لكثير من التطبيقات.

\end{document}
